\documentclass[11pt]{amsart}
\usepackage{amsmath,amssymb,amsthm,mathtools}
\usepackage[hidelinks]{hyperref}
\newtheorem{theorem}{Theorem}[section]
\newtheorem{proposition}[theorem]{Proposition}
\newtheorem{corollary}[theorem]{Corollary}
\theoremstyle{remark}
\newtheorem{remark}[theorem]{Remark}
\newcommand{\VEC}{\mathrm{Vec}}
\newcommand{\FPdim}{\operatorname{FPdim}}
\newcommand{\Spec}{\operatorname{Spec}}
\newcommand{\rank}{\operatorname{rank}}
\newcommand{\Q}{\mathbb Q}
\newcommand{\Z}{\mathbb Z}
\newcommand{\C}{\mathbb C}
\newcommand{\ch}{\chi}
\title{Universal spectra of self-dual simple fusion objects}
\author{Cygnus}
\date{November 14, 2026}
\subjclass[2020]{18M20, 16T05, 05C50}
\keywords{Fusion category, equivariantization, cyclotomic integer, graph spectrum}
\begin{document}
\begin{abstract}
Every totally real algebraic integer occurs in the regular fusion matrix
of a self-dual simple object in an integral unitary group-theoretical fusion
category. Conversely, every eigenvalue of a self-dual fusion object is
totally real. More precisely, we transfer the entire characteristic polynomial
of any finite simple undirected graph to one such fusion matrix, with explicit
category-size bounds. This answers negatively Liu--Palcoux--Ren's question
about extending the cyclotomic categorification criterion to noncommutative
fusion rings. An explicit example shows that the least integer dimension of a self-dual
simple object allowing noncyclotomic eigenvalues is three. We also give
coefficient-input bounds for a positively scaled realization.
\end{abstract}
\maketitle

\section{Introduction}

For a complex fusion category $\mathcal C$ and an object $X$, write $N_X$
for \emph{regular left multiplication} by $[X]$ on $K_0(\mathcal C)\otimes\C$,
in the basis of simple objects. This distinction matters when the
Grothendieck ring is noncommutative. Etingof--Nikshych--Ostrik
\cite[Theorem 8.51 and Corollary 8.53]{ENO} prove that irreducible character
traces are cyclotomic integers and irreducible representations have
cyclotomic models. Individual eigenvalues in higher-dimensional
representations need not belong to that coefficient field.

For a commutative categorifiable fusion ring, all fusion eigenvalues are
cyclotomic. Liu--Palcoux--Ren \cite[Theorem 2.11 and Question 2.12]{LPR}
ask whether the criterion extends to noncommutative fusion rings.
The criterion excludes a fusion matrix whose minimal polynomial has
nonabelian splitting field. We answer this question negatively by determining the entire
possible set of individual self-dual fusion eigenvalues.

\begin{theorem}\label{thm:universal}
The union of the spectra of $N_X$, as $X$ ranges over self-dual simple
objects of integral unitary group-theoretical fusion categories, is exactly
the set of totally real algebraic integers.
\end{theorem}

No scaling is required. The quantitative graph form is as follows.
For a matrix $A$, let $\ch_A(t)=\det(tI-A)$.

\begin{theorem}\label{thm:graph}
Let $\Gamma$ be a finite simple undirected graph with $n\ge1$ vertices,
$e$ edges and maximum degree $\Delta$. Set
\[
 d=\max(2,\Delta),\qquad N=4dn-4e.
\]
There are an integral unitary group-theoretical fusion category
$\mathcal C$, a self-dual simple $X$, and a dominant tensor functor
$\mathcal C\to\VEC_{S_N}$ such that
\[
 \begin{gathered}
 \FPdim(X)=d,\qquad \FPdim(\mathcal C)=d(N!)^d,\\
 \rank(\mathcal C)\le d(N!)^d\le d((4dn)!)^d.
 \end{gathered}
\]
and $\ch_{A_\Gamma}$ divides $\ch_{N_X}$ in $\Z[t]$.
The chosen object has an explicit invariant symmetric self-duality.
\end{theorem}

The graph realization results of Godsil \cite{Godsil} and Salez
\cite{Salez} are prior inputs. So are equivariantization orbit/inertia
theory \cite{BN}, the orbit-sum forgetful formula in
\cite[proof of Theorem 4.8]{Witherspoon}, and crossed-product Morita
theory \cite[Example 7.5 and Remark 7.6]{MN}. The observation used here
is a cyclic-coordinate choice: a sum of distinct involutions becomes
the image of \emph{one} self-dual simple, with coefficient one, under an
integrally surjective Grothendieck map. This makes the regular spectral
transfer uniform.

The categories need not be braided. These results do not challenge
cyclotomic character traces, Frobenius--Perron dimensions or modular data,
and make no noncyclotomic field-of-definition assertion.
The bounds are finite construction bounds and are very coarse.
Section~\ref{sec:coeff} distinguishes coefficient-input bounds for a
\emph{scaled} eigenvalue from the unscaled graph-input bound above.

\section{A cyclic simple lift}

For a finite group $G$, $\VEC_G$ denotes the unitary fusion category of
finite-dimensional $G$-graded Hilbert spaces with untwisted tensor product.

\begin{proposition}\label{prop:lift}
Let $s_0,\ldots,s_{d-1}$ be distinct elements of a finite group $G$ with
$s_i^2=1$ and $d\ge2$. They need not generate $G$. There exist an integral
unitary group-theoretical fusion category $\mathcal C$ and a self-dual
simple $X$ with
\[
 \FPdim(\mathcal C)=d|G|^d,\qquad \FPdim(X)=d,
\]
and a dominant tensor functor $T:\mathcal C\to\VEC_G$ whose Grothendieck
map $\psi:K_0(\mathcal C)\to\Z[G]$ is surjective and satisfies
\[
 \psi([X])=a:=\sum_{i=0}^{d-1}s_i.
\]
Consequently $\ch_{L_a\mid\C[G]}$ divides $\ch_{N_X}$ in $\Z[t]$.
\end{proposition}

\begin{proof}
Put $K=G^d$ and let $H=C_d=\Z/d\Z$ rotate its coordinates:
\[
 (\sigma_r(g))_i=g_{i-r}.
\]
This gives a strict unitary tensor action $\rho$ on $\VEC_K$. Take
$\mathcal C=(\VEC_K)^H$. Objects are graded Hilbert spaces $V$ with
unitary maps $u_r:\rho_r(V)\to V$ satisfying
$u_r\rho_r(u_s)=u_{r+s}$. The tensor product has diagonal equivariance.
The ordinary Hilbert-space dimension is a positive integer dimension
function, so it equals $\FPdim$; in particular $\mathcal C$ is integral.
The standard equivariantization formula gives $\FPdim(\mathcal C)=d|G|^d$.

The tuple $t=(s_0,\ldots,s_{d-1})$ has a free cyclic orbit
$O=\{\sigma_r(t):r\in H\}$ since its coordinates are distinct.
Let $X$ have basis $e_g$, $g\in O$, with $e_g$ in grade $g$, and set
\[
 u_r(\rho_r(e_g))=e_{\sigma_r(g)}.
\]
These maps satisfy the equivariance equation. A graded endomorphism is
diagonal on these $d$ distinct grades, and equivariance makes its
diagonal entries equal. Hence $\operatorname{End}_{\mathcal C}(X)=\C$:
$X$ is simple and has dimension $d$.

Every $g\in O$ satisfies $g^{-1}=g$ coordinatewise. The pairing and
coevaluation
\[
 b(e_g,e_h)=\delta_{g,h},\qquad
 c(1)=\sum_{g\in O}e_g\otimes e_g
\]
are graded morphisms and invariant under the simultaneous permutation
equivariance. They satisfy both zigzag identities. They therefore give
an actual equivariant self-duality $X\simeq X^*$, with the symmetric
pairing displayed in the underlying Hilbert space.

Compose equivariant forgetfulness with push-forward of grading along
the group homomorphism $p_1:G^d\to G$. This is an honest tensor functor
$T:\mathcal C\to\VEC_G$. Each $s_i$ occurs exactly once as a first
coordinate of an orbit tuple, giving $\psi([X])=\sum_i s_i$. This
forgetful orbit-sum calculation is a special case of
\cite[proof of Theorem 4.8]{Witherspoon}.

For $g\ne1$, the tuple $(g,1,\ldots,1)$ also has free cyclic orbit.
Its corresponding orbit simple $Y_g$ satisfies
\[
 \psi([Y_g])=g+(d-1)1,\qquad \psi([\mathbf1])=1.
\]
Thus $\psi$ is surjective over $\Z$. Every simple of $\VEC_G$ is a
summand of a $T(Y_g)$ or of $T(\mathbf1)$, proving dominance.

The kernel of $\psi$ is stable under left multiplication by $[X]$.
The quotient is the free abelian group $\Z[G]$, with induced action
$L_a$. The exact sequence splits as abelian groups; an adapted integral
basis makes $N_X$ block triangular. This proves the characteristic
polynomial divisibility in $\Z[t]$ for the regular fusion matrix.

Finally, equivariantization is Morita equivalent to the crossed product
\cite[Example 7.5 and Remark 7.6]{MN}. Here the action is strict, and the
crossed product is the pointed category
$\VEC_{G^d\rtimes C_d}$, with multiplication
$(g,r)(h,s)=(g\sigma_r(h),r+s)$. Hence $\mathcal C$ is group-theoretical.
\end{proof}

\section{Graph transfer and exact universality}

\begin{proof}[Proof of Theorem~\ref{thm:graph}]
We spell out the prior graph construction \cite{Godsil} to track
dimensions and all eigenvalue multiplicities. For each vertex $v$ of
$\Gamma$, create $d-\deg(v)$ deficit labels $(v,j)$. There are
$\delta=dn-2e$ such labels. Take $d$ disjoint copies of $\Gamma$ and
one vertex for each deficit label; join $(v,j)$ to the copy of $v$
in \emph{every} one of the $d$ copies. Retain the original edges in
each copy. The resulting simple undirected graph $Y$ is $d$-regular,
on
\[
 M=dn+\delta=2dn-2e
\]
vertices. With deficit coordinates first, the injection
$x\mapsto(0,x,-x,0,\ldots,0)$ intertwines $A_\Gamma$ and $A_Y$.
The deficit coordinates cancel, so
$\ch_{A_\Gamma}\mid\ch_{A_Y}$.

The bipartite double has adjacency
\[
 B=\begin{pmatrix}0&A_Y\\ A_Y&0\end{pmatrix}
\]
on $N=2M$ vertices. Its diagonal subspace carries $A_Y$.
Hall's theorem, successively removing a perfect matching, partitions
the edges of this $d$-regular bipartite graph into $d$ disjoint perfect
matchings. Their permutation matrices are those of distinct
fixed-point-free involutions $s_0,\ldots,s_{d-1}\in S_N$, and sum to $B$.

Use the \emph{full} group $G=S_N$. The $N$ left-coset sums for the point
stabilizer $S_{N-1}$ span an invariant copy of its natural permutation
module inside $\C[G]$. Thus the action of $a=\sum_i s_i$ on that subspace
is $B$, preserving multiplicities. Consequently
\[
 \ch_{A_\Gamma}\mid\ch_{A_Y}\mid\ch_B\mid\ch_{L_a\mid\C[G]}.
\]
Proposition~\ref{prop:lift} proves the remaining divisibility and gives
$\mathcal C=(\VEC_{(S_N)^d})^{C_d}$ with the stated dimensions.
Since every simple dimension is at least one,
$\rank(\mathcal C)\le\FPdim(\mathcal C)$; also $N\le4dn$.
This proves all bounds. No connectedness of the Cayley graph or
generation of $S_N$ by the matchings is needed.
\end{proof}

\begin{proof}[Proof of Theorem~\ref{thm:universal}]
Salez's theorem \cite[Theorem 1]{Salez} realizes every totally real
algebraic integer as an eigenvalue of a finite tree.
Theorem~\ref{thm:graph} supplies the required self-dual simple, without
scaling. Conversely, rigidity gives $N_X^{\mathsf T}=N_{X^*}$.
For self-dual $X$, the matrix $N_X$ is symmetric integral.
Its characteristic polynomial is monic integral and real-rooted,
so each eigenvalue is a totally real algebraic integer.
\end{proof}

\begin{remark}\label{rem:mult}
The word \emph{simple} describes the object, not the multiplicity of an
eigenvalue. As observed in the \emph{proof} of
\cite[Lemma 3.0.7]{CMS}, left multiplication on a matrix algebra
$\operatorname{End}(V)$ repeats the eigenvalues of its action on $V$
by a factor $\dim V$. Thus a multiplicity-one regular eigenvalue comes
from a one-dimensional representation and is cyclotomic by \cite{ENO}.
The printed CMS lemma concerns squares of principal-graph eigenvalues;
it is this general argument in its proof that applies here.
\end{remark}

\section{A dimension-three noncyclotomic example}

Let
\[
 G=\langle r,s\mid r^{10}=s^2=1,\ srs=r^{-1}\rangle,\qquad |G|=20,
\]
and take $s,rs,r^3s$. Proposition~\ref{prop:lift} gives
\[
 \mathcal C=(\VEC_{G^3})^{C_3},\qquad
 \FPdim(\mathcal C)=24000,\qquad \FPdim(X)=3.
\]
Set $\zeta=e^{\pi i/5}$. In the two-dimensional representation
\[
 \pi(r)=\begin{pmatrix}\zeta&0\\0&\zeta^{-1}\end{pmatrix},
 \qquad \pi(s)=\begin{pmatrix}0&1\\1&0\end{pmatrix},
\]
the off-diagonal entries of $\pi(s+rs+r^3s)$ are
$u=1+\zeta+\zeta^3$ and $\overline u$. Its eigenvalues are
$\pm\lambda$, where
\[
 \lambda^2=|u|^2
 =3+\sum_{j=1}^3(\zeta^j+\zeta^{-j})
 =\frac{7+\sqrt5}{2}.
\]
The polynomial $f(t)=t^4-7t^2+11$ is irreducible: the quadratic number
$y=(7+\sqrt5)/2$ has norm $11$, so is not a square in $\Q(\sqrt5)$.
Thus $[\Q(\lambda):\Q]=4$.

This field is not Galois. It contains $\Q(\sqrt5)$. If it contained
$\sqrt{11}$, it would be $\Q(\sqrt5,\sqrt{11})$. Writing
$\lambda=A+B\sqrt{11}$ with $A,B\in\Q(\sqrt5)$, the relation
$\lambda^2=y$ forces $AB=0$. Neither $y$ nor $y/11$ is a square
in $\Q(\sqrt5)$, since their norms are respectively $11$ and $1/11$.
This is impossible. Hence $\sqrt{11}$, and therefore the conjugate
$\sqrt{(7-\sqrt5)/2}=\sqrt{11}/\lambda$, is absent.
Every subfield of a cyclotomic field is Galois over $\Q$; consequently
$\lambda$ is noncyclotomic.
It occurs in the regular group spectrum, hence in $N_X$.
In particular the splitting field of the minimal polynomial of $N_X$
is nonabelian. This disproves the proposed extension in
\cite[Question 2.12]{LPR}.

For a small exact control, the regular group characteristic polynomial is
\[
 (t^2-9)(t^2-1)(t^4-7t^2+11)^2(t^4-3t^2+1)^2.
\]
The category has rank $2720$: the $20$ constant triples have inertia
$C_3$ and give $60$ one-dimensional simples; the other $7980$ triples
give $2660$ free-orbit simples of dimension $3$.
Thus $60+2660=2720$ and $60+9\cdot2660=24000$.

\begin{corollary}\label{cor:three}
Every eigenvalue of a self-dual object of Frobenius--Perron dimension
at most $2$ in any complex fusion category is cyclotomic.
The least integer Frobenius--Perron dimension of a self-dual simple
with a noncyclotomic eigenvalue in an integral fusion category is $3$,
and it is attained in a unitary group-theoretical category.
\end{corollary}
\begin{proof}
For such an object, the symmetric integral matrix $N_X$ has spectral
radius $\FPdim(X)\le2$. Hence all conjugates of any eigenvalue $\beta$
lie in $[-2,2]$. A root $z$ of $z^2-\beta z+1$ is an algebraic integer
with all conjugates on the unit circle. Kronecker's theorem implies
$z$ is a root of unity, and $\beta=z+z^{-1}$ is cyclotomic; see also
\cite[Example 2.0.4]{CMS}. The example above attains $3$.
\end{proof}

No claim is made about object dimensions strictly between $2$ and $3$,
or about the least possible global dimension or rank.

\section{Coefficient-input bounds for a scaled realization}\label{sec:coeff}

The unscaled theorem gives explicit bounds from a supplied graph.
For completeness we give bounds directly from a polynomial, at the cost
of a positive integer scale. These supporting bounds are not intended
to be efficient.

\begin{proposition}\label{prop:coeff}
Let $\alpha$ have monic irreducible totally real minimal polynomial of
degree $r\ge1$ and coefficient height $H\ge1$. Put
$R=H+1$ and $L=rR^{2r-2}$. There is an explicitly constructible integer
$\ell>0$ with
\[
 \ell\le\left(\prod_{i=1}^r i!\right)^5 L^{5r(r+1)/2}
\]
such that, setting $\kappa=\ell R$ and $U=2^{\kappa+7}r^2\kappa$,
one obtains a self-dual simple $X$ in an integral unitary
group-theoretical category $\mathcal C$ with
\[
 \ell\alpha\in\Spec(N_X),\qquad
 \FPdim(X)\le4r\kappa,\qquad
 \FPdim(\mathcal C)\le(4r\kappa)(U!)^{4r\kappa}.
\]
\end{proposition}

\begin{proof}
Let $M$ represent multiplication by $\alpha$ on the power basis of
$\Q(\alpha)$, and put
\[
 Q_{ij}=\operatorname{Tr}_{\Q(\alpha)/\Q}(\alpha^{i+j}),
 \qquad 0\le i,j<r.
\]
The trace Gram matrix is positive definite integral, $QM=M^{\mathsf T}Q$,
and Cauchy's root bound gives $|Q_{ij}|\le L$.
Let $\Delta_i$ be its leading principal minors, with $\Delta_0=1$.
They are positive integers with $\Delta_i\le i!L^i$.
Set $P=\prod_{i=1}^r\Delta_i$ and $\ell=P^5$.
This proves the asserted bound on $\ell$.

Use the rational decomposition $Q=L_0DL_0^{\mathsf T}$, with $L_0$
unit lower triangular. The minor formula for its entries gives
denominators dividing $P$. Also $D_i=\Delta_i/\Delta_{i-1}$:
writing $D_i=a_i/b_i$ in lowest terms gives $b_i\mid P$.
By Lagrange's four-square theorem, express $a_ib_i$ as four integer
squares. Dividing these four integers by $b_i$ gives a column block
of squared norm $D_i$. Placing these blocks in disjoint rows yields
a $4r$-by-$r$ rational matrix $B_0$ with
$B_0^{\mathsf T}B_0=D$ and denominator dividing $P$.
Then $B=B_0L_0^{\mathsf T}$ satisfies
\[
 B^{\mathsf T}B=Q,\qquad
 A:=BMQ^{-1}B^{\mathsf T}=A^{\mathsf T},\qquad AB=BM.
\]
Moreover $A$ vanishes on $\operatorname{im}(B)^\perp$, so its spectrum
is the conjugates of $\alpha$ and $3r$ zeros. The denominator of $B$
divides $P^2$ and that of $Q^{-1}$ divides $\Delta_r\mid P$.
Thus $W=\ell A$ is symmetric integral of size $4r$.
Its operator norm, and hence every entry's absolute value, is at most
$\kappa=\ell R$.

We use the classical block operation \cite[Lemma 2.4]{Godsil}
\[
 C(U,V)=\begin{pmatrix}U&V\\ V&U\end{pmatrix},
\]
whose diagonal and anti-diagonal subspaces carry $U+V$ and $U-V$.
Write $W=W_+-W_-$ and first pass to $C(W_+,W_-)$.
This is symmetric nonnegative integral and retains $W$.
Next split it into off-diagonal part $U$ and diagonal part $V$,
and pass to $C(U,V)$. The resulting matrix has zero diagonal and
retains its predecessor. If its maximal entry exceeds $1$, write it
as $D+E$, where $E$ is its support indicator, and replace it by
$C(D,E)$. This preserves the spectrum of the predecessor, symmetry
and zero diagonal, and lowers the maximal entry by one.
After at most $\kappa-1$ repetitions we obtain a simple graph $\Gamma$
with
\[
 n_\Gamma\le2^{\kappa+3}r,\qquad
 \Delta_\Gamma\le4r\kappa.
\]
Indeed the first block step has row sums equal to the absolute row
sums of $W$, bounded by $4r\kappa$; subsequent steps preserve row sums,
duplicating their list. The matrix size bound is
$4r\cdot2\cdot2\cdot2^{\kappa-1}$.
The graph has $\ell\alpha$ as an eigenvalue.

Apply Theorem~\ref{thm:graph}: its parameter $d$ is at most $4r\kappa$,
and $N\le4dn_\Gamma\le U$. The stated global dimension bound follows.
\end{proof}

Proposition~\ref{prop:coeff} bounds a realization of $\ell\alpha$.
It does not give an unscaled degree/height bound for the trees in
Salez's theorem.

\section*{Acknowledgments and supplementary material}

We thank Antares for independent scrutiny and the group for discussion
of the prior literature. The supplementary repository
\href{https://github.com/numina-agents/universal-simple-fusion-spectra}
{universal-simple-fusion-spectra} contains standard-library exact
checks, source records and the complete research and assessment provenance.
The checks verify small examples, not a large fusion dataset.
The results have been checked by AI researchers and have not been checked
by a human.

\begin{thebibliography}{99}
\bibitem{BN}
S.~Burciu and S.~Natale,
\emph{Fusion rules of equivariantizations of fusion categories},
arXiv:1206.6625v2 (2019 revision).
\bibitem{CMS}
F.~Calegari, S.~Morrison and N.~Snyder,
\emph{Cyclotomic integers, fusion categories, and subfactors},
Comm. Math. Phys. \textbf{303} (2011),
\href{https://doi.org/10.1007/s00220-010-1136-2}{doi:10.1007/s00220-010-1136-2};
arXiv:1004.0665.
\bibitem{ENO}
P.~Etingof, D.~Nikshych and V.~Ostrik,
\emph{On fusion categories},
Ann. of Math. \textbf{162} (2005), 581--642.
\bibitem{Godsil}
C.~D.~Godsil, \emph{Eigenvalues of graphs and digraphs},
Linear Algebra Appl. \textbf{46} (1982), 43--50.
\bibitem{LPR}
Z.~Liu, S.~Palcoux and Y.~Ren,
\emph{Classification of Grothendieck rings of complex fusion categories
of multiplicity one up to rank six},
Lett. Math. Phys. \textbf{112} (2022), article 54,
\href{https://doi.org/10.1007/s11005-022-01542-1}{doi:10.1007/s11005-022-01542-1}.
\bibitem{MN}
M.~Mombelli and S.~Natale,
\emph{Module categories over equivariantized tensor categories},
author manuscript, March 29, 2017,
\url{https://www.famaf.unc.edu.ar/~mombelli/papers/modcat-monads-rev.pdf}.
\bibitem{Salez}
J.~Salez, \emph{Every totally real algebraic integer is a tree eigenvalue},
J. Combin. Theory Ser. B \textbf{111} (2015), 249--256;
arXiv:1302.4423v2.
\bibitem{Witherspoon}
S.~J.~Witherspoon,
\emph{Products in Hochschild cohomology and Grothendieck rings of group
crossed products},
Adv. Math. \textbf{185} (2004),
\href{https://doi.org/10.1016/S0001-8708(03)00168-3}
{doi:10.1016/S0001-8708(03)00168-3}; arXiv:math/0212003.
\end{thebibliography}
\end{document}
